% Solution_Template.tex
\documentclass[11pt,a4paper]{article}
\usepackage[utf8]{inputenc}
\usepackage[T1]{fontenc}
\usepackage[ngerman, english]{babel} % remove ngerman if you prefer English only
\usepackage{geometry}
\usepackage{fancyhdr}
\usepackage{hyperref}

\geometry{margin=25mm}
\pagestyle{fancy}
\fancyhf{}
\renewcommand{\headrulewidth}{0.4pt}
\lhead{\textbf{Name:} \studentname}
\rhead{\textbf{Matr.-Nr.:} \studentnumber}
\cfoot{\thepage}

% User info - edit these
\newcommand{\studentname}{Matthias Janßen}
\newcommand{\studentnumber}{1871808}
\newcommand{\course}{Übersetzung und Analyse von Programmen}
\newcommand{\institute}{Universität Trier}

\title{\course}
\author{\studentname\\Matrikelnummer: \studentnumber\\\institute}
\date{\today}

\begin{document}

\maketitle

\begin{center}
    \textbf{Abgabe:} Assignment / Sheet \\
    \vspace{4mm}
    \textbf{Betreuer:} Prof. Dr. Stephan Diehl\\Jan Niclas Ruppenthal, M.Sc.
\end{center}

\bigskip

\section*{Aufgaben}
% Start writing your solutions here.

\section{Aufgabe 1}
% ...
Annahme: b = 2\\
Statische Variablen / Ausdrücke: 

    b $\rightarrow$ 2

    b > 0\\
Dynamische Variablen / Ausdrücke:

a, z

a > 0

a = a - b - z\\\\

Loop Unfolding:
Da b statisch 2 ist kann die while Schleife entfaltet werden. Im Schleifenrumpf wird b durch 2 ersetzt. Nach dem ersten Durchlauf
wird b um 1 reduziert. Dadurch entsteht der Ausdruck a = a - 1 - z.
\begin{verbatim}
    while (a > 0) {
        a = a - 2 - z
    }
    b = 2 - 1
    while (a > 0) {
        a = a - 1 - z
    }
\end{verbatim}

Beide Schleifen haben die gleiche Abbruchbedingung $a > 0$. Nach Terminierung der ersten Schleife gilt $a \le 0$. Zwischen den beiden Schleifen wird $a$ nicht verändert, daher ist die Bedingung der zweiten Schleife sofort falsch und ihr Körper unerreichbar. Falls die erste Schleife nicht terminiert, wird die zweite Schleife ohnehin nie erreicht.\\
Zur Optimierung kann die zweite Schleife entfernt werden, da ihr Körper in keinem Fall ausgeführt wird.\\
Optimierte Version:
\begin{verbatim}
    while (a > 0) {
        a = a - 2 - z
    }
\end{verbatim}
\section{Aufgabe 2}
% ...
Annahme: a = 1\\
Statische Variablen / Ausdrücke: 

    a $\rightarrow$ 1

    fib(4) $\rightarrow$ 5\\

Dynamische Variablen / Ausdrücke:

    x, y ,z, n

    x > a

    z = y

    x = x - 1

    y = y + z

    n <= 1

    n - 1

    n - 2\\\\
Code mit a = 1 eingesetzt:
\begin{verbatim}
    let fib(n) {if (n <= 1) then 1 else fib(n-1) + fib(n-2)}
        f(x, y, z) {
            z = y
            while (x > 1) do {
                x = x - 1
                y = y + z
            }
        }
    in fib(4) + fib(f(3, z, 1)) + f(z, 3, 1)
\end{verbatim}
Der Methodenaufruf fib(4) kann ausgewertet werden, da alle Argumente statisch sind. Es gilt fib(4) = 5.\\

\begin{verbatim}
    let fib(n) {if (n <= 1) then 1 else fib(n-1) + fib(n-2)}
        f(x, y, z) {
            z = y
            while (x > 1) do {
                x = x - 1
                y = y + z
            }
        }
    in 5 + fib(f(3, z, 1)) + f(z, 3, 1)
\end{verbatim}
Polyvariante Spezialisierung von f:
Beim Aufruf f(3, z, 1) kann die Schleife entfaltet werden, da x und a statisch sind (2 Iterationen):

\begin{verbatim}
    f_x3(y) {
        z = y
        y = y + z
        y = y + z
    }
\end{verbatim}

Für f(z, 3, 1) ist das nicht möglich, da x dynamisch ist. y wird jedoch statisch auf 3 gesetzt:

\begin{verbatim}
    f_y3(x) {
        y = 3
        z = 3
        while (x > 1) do {
            x = x - 1
            y = y + 3
        }
    }
\end{verbatim}

Gesamter partiell ausgewerteter Code:

\begin{verbatim}
    let fib(n) {if (n <= 1) then 1 else fib(n-1) + fib(n-2)}
        f_x3(y) {
            z = y
            y = y + z
            y = y + z
        }

        f_y3(x) {
            y = 3
            z = 3
            while (x > 1) do {
                x = x - 1
                y = y + 3
            }
        }

    in 5 + fib(f_x3(z)) + f_y3(z)

\end{verbatim}

\end{document}